\mysec{Lista 4}
\etocsetnexttocdepth{subsection}
\localtableofcontents
\newcounter{quatro}
\stepcounter{quatro}

%ex 1
\myexe{\thequatro}
\bigline
\textbf{Seja $(X, \mcal F, \mu)$ um espaço de medida e $f \in L^p(X, \mu)$, 
com $1 \le p < +\infty$. Prove a Desigualdade de Chebyshev:
\begin{align*}
  \mu\parth{\set{x \in X \st |f(x)| > \alpha}} \le \cparth{\frac{\pnorm p f}{\alpha}}^p, \quad \forall \alpha>0.
\end{align*}}
\begin{proof}[\textbf{Dem:}] Como $f \in L^p (X, \mu)$, temos $|f|^p \in L^1(X, \mu)$, ou seja, 
  \begin{align*}
    \int |f|^p d\mu < + \infty.
  \end{align*}
  Mais ainda, como $|f|^p \ge 0$, temos que em particular que todos os representantes da classe 
  $[|f|^p]$ são mensuráveis não-negativos. Portanto, 
  da desigualdade de Chebyshev para funções em $\fplus M (X, \mcal F)$, dados $\beta = \alpha^p > 0$ e
  $E_\beta = \set{x \in X \st |f(x)|^p > \beta}$,
  \begin{align*}
    \mu(\set{x \in X \st |f(x)| > \alpha}) = \mu(E_\beta) \le \frac{1}{\beta} \int |f|^p d\mu 
    = \frac{1}{\alpha^p} \pnorm p f^p.
  \end{align*}
\end{proof}
\bigline

%ex 2 
\stepcounter{quatro}
\myexe{\thequatro}

\noindent 
\textbf{Sejam $(X, \mcal F, \mu)$ um espaço de medida e $f \in L^p (X, \mu)$, com $1 \le p < +\infty$. 
Mostre que o conjunto $E = \set{x \in X \st |f(x)| \neq 0}$ é $\sigma$-finito.}
\begin{proof}[\textbf{Dem:}] Note que $|f(x)|^p \ge 0$ para cada 
  $x \in X$ e $1 \le p \le +\infty$. Para $p \neq +\infty$, considere  
  $E_n = \set{x \in X \st |f(x)| > \frac{1}{n}}$. 
  Note que $E = \bigcup_{n \in \bN} E_n$, desse modo, vejamos que 
  $E_n$ é finito. Note que da desigualde de Chebyshev, para 
  cada $n \in \bN$,
  \begin{align*}
    \mu(E_n) \le \cparth{\frac{||f||_p}{n}}^p < +\infty.
  \end{align*}
  Com isso, segue que $E$ é $\sigma$-finito.
\end{proof}

%ex 3
\nextexe{quatro}
\textbf{Sejam $(X, \mcal F, \mu)$ um espaço de medida e $f \in L^p(X, \mu)$, 
com $1 \le p \le +\infty$. Definindo $E_n = \set{x \in X \st |f(x)| \ge n}$, 
mostre que $\lim \mu(E_n) = 0$.}
\begin{proof}[\textbf{Dem:}] Note que, para cada $n \in \bN$ temos 
$E_n \supseteq E_{n+1}$. Mais ainda, se $p \neq +\infty$, para cada $n \in \bN$, 
  da desigualdade de Chebyshev para espaços $L^p(X, \mu)$, dado $f \in L^p(X, \mu)$, 
\begin{align*}
  \mu(E_n) \le \cparth{\frac{||f||_p}{n}}^p \implies 
  \lim_{n \rightarrow \infty} \mu(E_n) \le 
  \lim_{n \rightarrow \infty} \cparth{\frac{||f||_p}{n}}^p = 0. 
\end{align*}
Se $p = +\infty$, por definição, dado $f \in L^\infty(X, \mu)$, 
\begin{align*}
  ||f||_\infty = \inf \set{a \ge 0 \st \mu(\set{x \in X \st |f(x)| > a}) = 0},
\end{align*}
  isto é, $\mu(\set{x \in X \st |f(x)| > ||f||_\infty}) = 0$.  
  Visto que $f \in L^\infty(X, \mu)$, existe $M \in \bN$ tal que 
$||f||_\infty < M$. Desse modo, para todo $n \ge M$,
\begin{align*}
  \mu(\set{x \in X \st |f(x)| \ge n}) \le 
  \mu\parth{\set{x \in X \st |f(x)| > ||f||_\infty}} = 0,
\end{align*}
  portanto, $\lim \mu(E_n) = 0$.
\end{proof}

%ex 4
\newexe{quatro}
\textbf{Sejam $X = \bN$ e $\mu$ a medida da contagem em $\bN$. 
Considere $f: \bN \rightarrow \bR$ definida por $f(n)=1/n$. Mostre que 
$f$ não pertence a $L^1(\bN, \mu)$, mas pertence a $L^p(\bN, \mu)$ 
para todo $1 < p \le \infty$.}
\begin{proof}[\textbf{Dem:}] Ora, como visto em exercícios anteriores, 
\begin{align*}
  \int |f| d\mu = \sum_{k \in \bN} |f(k)| = \sum_{k \in \bN} \mparth {\frac{1}{k}} 
  = \sum_{k \in \bN} \frac{1}{k},
\end{align*}
com isso, segue que $\int |f| d\mu$ é igual a série harmônica, i.e., $\int |f| d\mu = +\infty$,
  ou seja, $f \not \in L^1(X, \mu)$. 
  Entretanto, 
\begin{align*}
  \int |f|^p d\mu = \sum_{k \in \bN} |f(k)|^p = \sum_{k \in \bN} \mparth{\frac{1}{k}}^p 
  = \sum_{k \in \bN} \frac{1}{k}^p,
\end{align*}
ou seja, da convergência das $p$-séries, segue que $\int |f|^p d\mu < +\infty$, 
  isto é, $f \in L^p(X, \mu)$ para 
  $1 < p < +\infty$. Agora, para $p = +\infty$, note que, para cada $n \in \bN$, 
\begin{align*}
  f(n) \le \max f(k) = 1 \implies ||f||_\infty \le \max f(k) = 1 < +\infty. 
\end{align*}
  Portanto $f \in L^\infty(X, \mu)$.
\end{proof}

%ex 5
\nextexe{quatro}
\textbf{Sejam $X = \bN$ e $\mu$ a medida da contagem em $\bN$. 
Prove que se $f \in L^p(\bN, \mu)$, então $f \in L^q(\bN, \mu)$ para todo $1 \le p \le q < \infty$, 
e $||f||_q \le ||f||_p$.}
\begin{proof}[\textbf{Dem:}] Note que para $p = q$ segue diretamente o resultado. 
  Desse modo, suponhamos que $f \in L^p(\bN, \mu)$. Já que $\mu$ é a medida da contagem 
  em $\bN$, 
  \begin{align*}
    ||f||_p = \parth{\int |f|^p d\mu}^{1/p} = \parth{\sum_{k \in \bN} |f(k)|^p}^{1/p}
    \ge |f(n)|, \quad \forall n \in \bN.
  \end{align*}
  Daqui, note que, para cada $n \in \bN$, dado $q-p > 0$, 
  \begin{align*}
    |f(n)|^q = |f(n)|^p\cdot |f(n)|^{q-p} \le |f(n)|^p \cdot ||f||_p^{q-p}
  \end{align*}
  Ou seja, 
  \begin{align*}
    \parth{\int |f|^q d\mu}^{1/q} \le \parth{\int |f|^p\cdot ||f||_p^{q-p} d\mu}^{1/q}
    = \parth{||f||_p^{q-p}\int |f|^pd\mu}^{1/q} = ||f||_p,
  \end{align*}
  isto é, $||f||_q \le ||f||_p$. Mais ainda, segue diretamente da desigualdade 
  que $f \in L^q(\bN, \mu)$.

\end{proof}

%ex 6
\nextexe{quatro}
\textbf{Seja $X = \bN$ e defina $\lambda: \mcal P(\mbb N) \rightarrow \bN$ por 
\begin{align*}
  \lambda(E) = \sum_{n \in E}\frac{1}{n^2}
\end{align*}
Considere $f: \bN \rightarrow \bR$ dada por $f(n)= \sqrt{n}$.
\begin{enumerate}[label=\alph*)]
  \item Mostre que $\lambda$ é uma medida em $\bN$ e que $\lambda(\bN) < +\infty$.
    \begin{proof}[\textbf{Dem:}]
      Por definição segue que $\lambda(E) \ge 0$ para todo $E \in \mcal P(\bN)$ não 
      vazio. Para $E = \emptyset$,defina $\lambda(E) = 0$. Daqui, vejamos que vale 
      a $\sigma$-aditividade. Dado $\set{E_j}$ uma sequência em $\mcal F$ dois a 
      dois disjunto, note que, 
      \begin{align*}
        \lambda \parth{\bigcup_{j \in \bN} E_j} = \sum_{n \in \bigcup E_j} \frac{1}{n^2}
        = \sum_{j \in \bN} \cparth{\sum_{n \in E_j} \frac{1}{n^2}} 
        = \sum_{j \in \bN} \lambda(E_j), 
      \end{align*}
      ou seja, $\lambda$ é uma medida em $\bN$. Como $\lambda(\bN)$ é a $p$-série 
      com $p=2$, segue que $\lambda(\bN) < +\infty$.
    \end{proof}
  \item Mostre que $f \in L^p(\bN, \lambda)$ se, e somente se, $1\le p < 2$.
  \begin{proof}[\textbf{Dem:}] Seja $p \ge 1$. Dado qualquer 
    $f: \bN \rightarrow \bR$ segue que para todo $E \subset \bR$, 
    $f^{-1}(E) \in \mathcal P (\bN)$, isto é, toda função $f: \bN \rightarrow \bR$
    é mensurável em $(X, \mathcal P(\bN))$. Daqui, note que dos exercícios anteriores,
    dado $|f|^p \in M(\bN, \mathcal P(\bN))$,
    \begin{align*}
      \int |f|^p d\lambda &= \sum_{n \in \bN} |f(n)|^p \lambda(\set{n})
      = \sum_{n \in \bN}  \cparth{|f(n)|^p \sum_{k \in \set{n}} \frac{1}{k^2}}\\
      &= \sum_{n \in \bN} \frac{|f(n)|^p}{n^2} = \sum_{n \in \bN} \frac{n^{p/2}}{n^2}
      = \sum_{n \in \bN} \frac{1}{n^{2-p/2}}.
    \end{align*}
    Da convergência das p-séries, segue que $\sum_{n \in \bN} \frac{1}{n^p}$ converge
    para algum número real se, 
    e somente se $p>1$. Desse modo, $\int |f|^p d\lambda < + \infty$ se, e somente se 
    $2-p/2 > 1$, isto é, $p < 2$. Portanto, $f \in L^P(\bN, \lambda)$ se, e somente se 
    $1\le p < 2$.
  \end{proof}
\end{enumerate}}

%ex 7
\nextexe{quatro}
\textbf{Seja $(X, \mcal F, \mu)$ um espaço de medida finita. Se $f$ é mensurável em $X$, 
defina $E_n = \set{x \in X \st (n - 1) \le |f(x)| < n}$. Mostre que $f \in L^1(X, \mu)$ 
se, e somente se, 
\begin{align*}
  \sum_{n \in \bN} n \mu(E_n) < +\infty.
\end{align*}
Generalize para $1 \le p < +\infty$, isto é, mostre que $f \in L^p(X, \mu)$, se 
e somente se 
\begin{align*}
  \sum_{n \in \bN} n^p \mu(E_n) < + \infty.
\end{align*}}
\begin{proof}[\textbf{Dem:}] Definindo $N = \{x \in X \st |f(x)| = +\infty \}$, segue que 
  $N^C = \{x \in X \st |f(x)| < +\infty\}$. E, como $N^C = \bigcup E_n$ onde para 
  cada $k, l \in \bN$, $E_k \cap E_l = \emptyset$ sempre que $k \neq l$, temos que 
  $\mu\parth{N^C} = \mu\parth{\bigcup E_n} = \sum \mu(E_n)$. Com isso, vejamos a 
  ida e a volta da proposição.\\
$(\implies)$ Suponhamos que $f \in L^1(X, \mu)$. Então, $|f| < +\infty$ $\mu$-q.t.p.. Mais 
  precisamente $\mu(N) = 0$. Daqui, note que,
\begin{align*}
  \int |f| d\mu &= \int_N |f| d\mu + \int_{N^C} |f| d\mu = \int_{N^C} |f| d\mu 
  = \int |f| \chi_{N^C} d\mu\\
  &= \int |f| \chi_{\bigcup E_n} d\mu = \int |f| \parth{\sum_{n \in \bN} \chi_{E_n}}d\mu < +\infty.
\end{align*}
Ou seja, 
\begin{align*}
  \int |f| d\mu = \int |f| \parth{\sum_{n \in \bN} \chi_{E_n}} = \sum_{n \in \bN} \int |f| \chi_{E_n} d\mu.
\end{align*}
Pela definição de $E_n$, temos,
\begin{align*}
  +\infty > \int |f| d\mu \ge \sum_{n \in \bN} \int (n-1) \chi_{E_n} d\mu
  = \sum_{n \in \bN} (n-1) \mu\parth{E_n}.
\end{align*}
  Como $N^C \subset X$ e $\mu(X) < +\infty$,
  \begin{align*}
    +\infty > \cparth{\sum_{n \in \bN} (n - 1) \mu\parth{E_n}} + 
    \cparth{\sum_{n \in \bN} \mu\parth{E_n}} 
    = \sum_{n \in \bN} n\mu\parth{E_n}.
  \end{align*}
  $(\impliedby)$ Como $f: X \rightarrow \bR$, note que $\mu({x \in X \st |f(x)| = +\infty}) = 0$. Portanto, 
  temos da definição de $f$ que $X = \bigcup E_n$. Já que $f$ é mensurável e limitada por $n$ em $E_n$
  para todo $n \in \bN$,
  \begin{align*}
    \int |f| d\mu = \int |f| \chi_X d\mu = \int |f| \chi_{\bigcup E_n} d\mu = \int \sum_{n \in \bN} |f| \chi_{E_n} d\mu 
    \le \int \sum_{n \in \bN} n \chi_{E_n} d\mu.
  \end{align*}
  Já que para cada $n \in \bN$, $n \chi_{E_n} \in \mathcal M^+(X, \mu)$,
  \begin{align*}
    \int |f| d\mu \le \int \sum_{n \in \bN} n \chi_{E_n} d\mu = \sum_{n \in \bN} \int n \chi_{E_n}d \mu, 
    = \sum_{n \in \bN} n \mu(E_n) < +\infty
  \end{align*}
  ou seja, $f \in L^1(X, \mu)$.
\end{proof}

%ex 8
\nextexe{quatro}
\textbf{Sejam $(X, \mcal F, \mu)$ um espaço de medida finita e $f \in L^p(X, \mcal F, \mu)$, $1 \le p < \infty$, 
e $\epsilon > 0$. Mostre que existe um conjunto $E_\epsilon \in \mcal F$ com $\mu(E_\epsilon) < +\infty$ 
tal que se $F \in \mcal F$ e $F \cap E_\epsilon = \emptyset$, então $\pnorm p {f \chi_F} < \epsilon$.}
\begin{proof}[\textbf{Dem:}] Ora, como $(X, \mcal F, \mu)$ é um espaço de medida finita, então, 
  para todo $A \in \mcal F$, segue que $\mu(A) < +\infty$. Agora, como $f \in L^p(X, \mcal F, \mu)$, 
  todo representante de $|f|^p \in L^1(X, \mu)$ é mensurável. Portanto, temos que 
  \begin{align*}
    E_\epsilon = \set{x \in X \st |f(x)|^p \ge \frac{\epsilon^p}{\mu(X)}} \in \mcal F.
  \end{align*}
  Agora, se $F \in \mcal F$ e $F \cap E_\epsilon = \emptyset$, segue que, 
  \begin{align*}
    F \subset \set{x \in X \st |f(x)|^p < \frac{\epsilon^p}{\mu(X)}}.
  \end{align*}
  E, como os representantes de $|f|^p$ são 
  mensuráveis, segue que os representantes de $|f|^p\chi_F = |f \chi_F|^p$ também são. Portanto, 
  \begin{align*}
    \int |f \chi_F|^p d\mu = \int | f |^p \chi_F d\mu < \int \frac{\epsilon^p}{\mu(X)} \chi_F d\mu 
    = \epsilon^p \frac{\int_F d\mu}{\mu(X)} = \epsilon^p \cparth{\frac{\mu(F)}{\mu(X)}} \le \epsilon^p.
  \end{align*}
Ou seja, 
\begin{align*}
  \pnorm p {f\chi_F} = \cparth{\int |f\chi_F|^p d\mu}^{1/p} < \epsilon.
\end{align*}
\end{proof}

%ex 9
\nextexe{quatro}
\textbf{Sejam $(X, \mcal F, \mu)$ um espaço de medida e $f \in L^\infty(X, \mu)$. 
Mostre que se $A < \pnorm \infty f$, então existe um conjunto 
$E \in \mcal F$ com $\mu(E) > 0$ tal que $|f(x)| > A$ para todo $x \in E$.}
\begin{proof}[\textbf{Dem:}] Por definição, 
  \begin{align*}
    \pnorm \infty f = \inf \set{\alpha \ge 0 \st \mu(\set{x \in X \st |f(x)| > \alpha}) = 0}
  \end{align*}  
  Desse modo, como $\alpha = ||f||_\infty$ é o menor real extendido satisfazendo 
  \begin{align*}
    \mu(\set{x \in X \st |f(x)| > \alpha}) = 0,
  \end{align*}
  já que $\pnorm \infty f > A$, então $\mu(\set{x \in X \st |f(x)| > A}) > 0$. Definindo 
  $E = \set{x \in X \st |f(x)| > A}$, segue que para todo $x \in E$, temos $|f(x)| > A$,
  mais ainda, o conjunto satisfaz $\mu(E) > 0$, como queríamos.
\end{proof}

%ex 10
\nextexe{quatro}
\textbf{Sejam $(X, \mcal F, \mu)$ um espaço de medida. Mostre que se $f \in L^P(X, \mu)$, $1 \le p \le \infty$, 
e $g \in L^\infty(X, \mu)$, então $fg \in L^p(X, \mu)$ e $\pnorm p {fg} \le \pnorm p f \pnorm \infty g$.}
\begin{proof}[\textbf{Dem:}] Para $p = \infty$, segue diretamente que, 
  \begin{align*}
    \pnorm \infty {fg} = \inf \set{\alpha \ge 0 \st \mu(\set{x \in X \st |f(x)g(x)| > \alpha}) = 0}.
  \end{align*}
  Entretanto, como 
  \begin{align*}
    \mu(\set{x \in X \st |f(x)| > \pnorm \infty f}) = 0, \quad \mu(\set{x \in X \st |g(x)| > \pnorm \infty g}) = 0, 
  \end{align*}
  então, 
  \begin{align*}
    \mu(\set{x \in X \st |f(x)| |g(x)| > \pnorm \infty f \pnorm \infty g}) 
    \le \mu(\set{x \in X \st |f(x)| > \pnorm \infty f}) = 0.
  \end{align*}
  Ou seja, da definição de infimo, segue que, 
  \begin{align*}
    \pnorm \infty {fg} \le \pnorm \infty f \pnorm \infty g.
  \end{align*}
  Já que $f, g \in L^\infty(X, \mu)$, temos da desigualdade acima que $fg \in L^\infty(X, \mu)$. \\ 
  Suponhamos agora que $ 1 \le p < \infty$, Note que, como os representantes 
  de $|f|^p \in L^p(X, \mu)$ são mensuráveis, segue que para todo $\alpha \in \bR$
  temos que $\alpha |f|^p$ é também mensurável.
  Portanto, como $g \in L^\infty(X, \mu)$, sabemos que $g$ é mensurável e 
  por conseguinte, que $|g|$ também é. Já que $\phi: \bR \rightarrow \bR$ definida 
  por $\phi(x) = x^p$ é, em particular para $1 \le p < \infty$, contínua, 
  segue que $|g|^p$ é mensurável, (ver exercício 10 Lista 1). Portanto, 
  $|fg|^p$ é mensurável, desse modo, já que $|fg|^p \le |f|^p\pnorm \infty g ^p$ $\mu$-q.t.p.,
  \begin{align*}
    \int |fg|^P d\mu &= \int |f|^p |g|^p d\mu \le \int |f|^p \pnorm \infty g^p d\mu \\ 
    &= \pnorm \infty g^p \int |f|^p d\mu,
  \end{align*}
  já que $|f|^p \in L^1(X, \mu)$, segue que $\int |fg|^p d\mu \le \pnorm \infty g^p 
  \pnorm p f^p < +\infty$, isto é, $|fg| \in L^p(X, \mu)$. Mais ainda, 
  como $\pnorm p {fg}^p = \int |fg|^p d\mu \le \cparth{\pnorm p f \pnorm \infty g}^p$,
  temos 
  $$\pnorm p {fg} \le \pnorm p f \pnorm \infty g.$$ 
  Como queríamos.
\end{proof}

\nextexe{quatro}
\textbf{Sejam $(X, \mcal F, \mu)$ um espaço de medida e $1 \le p < r \le +\infty$. 
Defina em $L^p(X, \mu) \cap L^r(X, \mu)$ a função $\pnorm {} {\,\cdot\,}:
L^p(X, \mu) \cap L^r(X, \mu) \rightarrow \bR$ dada por 
$\pnorm {} f \coloneq \pnorm p f + \pnorm r f$. Mostre que 
$\pnorm {} {\, \cdot \,}$ é uma norma e que $L^p(X,\mu)\cap L^r(X, \mu)$ é 
um Espaço de Banach com esta norma.}
\begin{proof}[\textbf{Dem:}]
Vejamos primeiro que $\pnorm {} {\, \cdot \,}$ é uma norma em $L^p(X, \mu) \cap L^r(X, \mu)$.
\begin{itemize}
  \item $\snorm f \ge 0$ e $\snorm{f} = 0 \iff f = 0$:\\
    Dado $f \in L^p(X, \mu) \cap L^r(X,\mu)$, como $\pnorm p {\, \cdot \,}$ e 
    $\pnorm r {\, \cdot \,}$ em $L^p(X, \mu) \cap L^r(X, \mu)$, segue que 
    $\snorm f = \pnorm p f + \pnorm r f \ge 0$. Agora, suponhamos que $\snorm f = 0$. Então, 
    \begin{align*}
      \snorm f = \pnorm p f + \pnorm r f = 0 \implies \pnorm p f = - \pnorm r f \implies 
      \pnorm p f = 0 \implies f = 0.
    \end{align*}
    Por outro lado, suponhamos que $f = 0$, desse modo, 
    \begin{align*}
      f = 0 \implies \pnorm p f = \pnorm r f = 0 \implies 
      \snorm f = \pnorm p f + \pnorm r f =0.
    \end{align*}
  \item Dados $\lambda \in \bR$ e $f \in L^p(X, \mu)\cap L^r(X, \mu)$,
    $\snorm{\lambda f} = \mparth{\lambda} \snorm f$:

    Ora, como $\pnorm p {\, \cdot \,}$ e 
    $\pnorm r {\, \cdot \,}$ são normas em $L^p(X, \mu) \cap L^r(X, \mu)$, 
    \begin{align*}
      \mparth{\mparth{\lambda f}} = \pnorm p {\lambda f} + \pnorm r {\lambda f}
      = \mparth \lambda \pnorm p f + \mparth \lambda \pnorm r f = \mparth \lambda \snorm f.
    \end{align*}
  \item Dados $f, g \in L^p(X,\mu)\cap L^r(X, \mu)$, $\snorm {f + g} \le \snorm f + \snorm g$:
    
    Novamente, como $\pnorm p {\, \cdot \,}$ e 
    $\pnorm r {\, \cdot \,}$ são normas em $L^p(X, \mu) \cap L^r(X, \mu)$, 
    \begin{align*}
      \snorm {f+g} &= \pnorm p {f + g} + \pnorm r {f + g} \le 
      \pnorm p f + \pnorm p g + \pnorm r f + \pnorm r g \\ 
      &= \parth{\pnorm p f + \pnorm r f} + \parth{\pnorm p g + \pnorm r g}
      = \snorm f + \snorm g. 
    \end{align*}
\end{itemize}
  Com isso, segue que $\snorm {\, \cdot \,}$ é norma em $L^p(X, \mu) \cap L^r(X, \mu)$.
  Daqui, vejamos que é também um Espaço de Banach quando munido da norma $\snorm {\, 
  \cdot \,}$. Para tal, vejamos que é completo, isto é, toda sequência 
  de Cauchy em $(L^p(X, \mu) \cap L^r(X, \mu), \snorm {\, \cdot \,})$ converge 
  no sentido $\snorm {\, \cdot \,}$ para alguém em $L^p(X, \mu) \cap L^r(X, \mu)$.
  Considere $(f_n)_{n \in \bN}$ uma sequência de Cauchy em $(L^p(X, \mu) \cap L^r(X, \mu), \snorm {\, \cdot \,})$, 
  ou seja, para todo $\epsilon > 0$, existe $N \in \bN$ tal que dados inteiros $n, m \ge N$, temos,
  \begin{align*}
    \snorm{f_n - f_m} < \epsilon.
  \end{align*}
  Ora, por definição, 
  \begin{align*}
    \pnorm p {f_n - f_m} &\le \pnorm p {f_n - f_m} + \pnorm r {f_n - f_m}  = \snorm {f_n - f_m} < \epsilon, \text{ e também, } \\
    \pnorm r {f_n - f_m} &\le \pnorm p {f_n - f_m}  + \pnorm r {f_n - f_m} = \snorm {f_n - f_m} < \epsilon.
  \end{align*}
  ou seja, $(f_n)_{n \in \bN}$ é Cauchy tanto em $(L^p(X, \mu), \pnorm p {\, \cdot \,})$, quanto 
  em $(L^r(X, \mu), \pnorm r {\, \cdot \,})$. Como ambos são Espaços de Banach, temos que, 
  $(f_n)_{n \in \bN}$ converge para $g \in L^p(X, \mu)$ e $h \in L^r(X, \mu)$  
  simultaneamente. Por modos de convergência,
  sabemos que, $(f_n)$ converge em medida 
  para $g$ e $h$ simultaneamente. Das propriedades de convergência em medida, 
  segue que $g = h$ $\mu$-q.t.p.. Desse modo, $g = h = f \in L^p(X, \mu)\cap L^r(X, \mu)$.
  Daqui, note que, para todo $\epsilon > 0$ existe $N \in \bN$ tal que 
  $\pnorm p {f_n - f} < \epsilon/2$ e $\pnorm r {f_n - f}< \epsilon /2 $, ou seja,
  \begin{align*}
    \snorm {f_n - f} = \pnorm p {f_n - f} + \pnorm r {f_n - f} < \epsilon,
  \end{align*}
  isto é, $(f_n)$ converge no sentido $\snorm {\, \cdot \,}$ para 
  $f \in L^p(X, \mu) \cap L^r(X, \mu)$, i.e., 
  $(L^p(X, \mu)\cap L^r(X, \mu), \snorm {\, \cdot \,})$ 
  é um Espaço Completo.
\end{proof}

%ex 12 
\nextexe{quatro}
\textbf{Sejam $(X, \mcal F, \mu)$ um espaço de medida e $1 \le p < r \le +\infty$. 
Defina em $L^p(X, \mu) + L^r(X, \mu)$ a função 
$\snorm {\, \cdot \,}: L^p(X, \mu) + L^r(X, \mu) \rightarrow \bR$ dada por 
\begin{align*}
  \snorm f \coloneq \inf \set{\pnorm p g + \pnorm r g \st 
  g \in L^p(X, \mu), h \in L^r (X, \mu), \text{ e } f = g+ h}
\end{align*}
Mostre que $\snorm {\, \cdot \,}$ é uma norma e 
que $L^p(X, \mu) + L^r(X, \mu)$ é um Espaço de Banach com esta norma.}
\begin{proof}[\textbf{Dem:}]
  
\end{proof}
